\chapter*{Executive Summary}
\addcontentsline{toc}{chapter}{Executive Summary}

\paragraph{Purpose and evidence basis.} This report documents an MSc
technical activity investigating how much can be learned about a
ferrocement specimen's condition from repeated corrosion photography. It
builds on a predecessor experimental dataset — 48 specimens fabricated,
chloride-exposed, and terminally tested across two campaigns, external to
the present activity — and works from the resulting photographic record,
specimen metadata, and terminal mechanical measurements. Visible-corrosion
labels are available for nearly every image in the current audited basis
(791 of 792 raw files), while wire-area loss and ultimate load are each
obtained only once per specimen, at the end of its exposure period. This
asymmetry between dense surface supervision and sparse terminal structural
supervision governs how every finding in this report should be read.

\paragraph{Technical activity.} The activity comprised: an audit and
alignment of the raw photographic and workbook data; three successive
approaches to visible-corrosion estimation (a classical handcrafted
baseline, a frozen deep-image embedding, and an interpretable
multi-family feature representation); a feasibility study of hidden
structural-condition inference; systematic robustness evaluation under
grouped, treatment-held-out, and campaign-held-out specimen partitions; a
refocused study explicitly testing whether images add value beyond
metadata for terminal ultimate load; an exploratory degradation-curve and
proxy-remaining-useful-life screening pipeline; preparation of a
leakage-safe dataset and augmentation pipeline for a four-class
corrosion-severity classification task; and an audit of the project's own
software and reproducibility practices.

\paragraph{Main findings.} The evidence supports a clear hierarchy of
confidence:
\begin{enumerate}
\item \textbf{Visible corrosion} is the strongest and most consistently
learnable image-based quantity in this dataset, across three distinct
representations.
\item \textbf{Hidden structural inference} is substantially weaker: the
available structural models are either mixed-input or, in their final,
most carefully selected form, metadata-only, and none transfers reliably
across campaigns.
\item \textbf{Terminal ultimate load}, one of the two directly measured
terminal structural endpoints and the endpoint selected for the
refocused analysis, is dominated by specimen and exposure metadata under
pooled evaluation; no robust incremental benefit from image-derived
features was established.
\item \textbf{Degradation and proxy-remaining-useful-life screening}
remain exploratory: the pipeline is a working demonstration of a
derivation chain, not a validated prognostic method.
\end{enumerate}
Two cross-cutting points qualify all four: which evaluation regime is used
materially changes what can be claimed, and leave-one-campaign-out is a
severe extrapolation stress test in this dataset because campaign identity
is entangled with steel-mesh configuration, chloride level, and exposure
duration. In addition, some strong visible-corrosion benchmarks rely on
image-derived features that are very closely aligned with the target,
creating shortcut or label-adjacency risk; the strongest case is the
near-reconstructive total-rust pairing identified by the later feature
audit, and this is checked and flagged explicitly where it applies.

\paragraph{Scope and limitations.} Every structural or prognostic
conclusion in this report is built on 48 terminal structural
measurements, from only two experimental campaigns whose design factors
are almost perfectly confounded with campaign identity. No specimen has a
repeated or intermediate structural measurement, and no specimen has an
observed failure or lifetime event; structural values reported at
non-terminal weeks are always model estimates, never measurements. The
four-class classification branch's data preparation is complete, but no
classifier has been trained, compared, or evaluated at the time of
writing.

\paragraph{Current status and next direction.} The main report body
(dataset audit, methodological framework, activity chronology, integrated
results, interpretation, software/reproducibility audit, limitations, and
conclusions) is complete. Reproducibility has improved over the course of
the project — deterministic augmentation, saved split manifests, and an
automated feature-diagnostics report are already in place — though two
live configuration and source-code issues, identified and located in this
report, still prevent an unmodified rerun of some pipeline stages. The
most immediately justified next modelling step is training a baseline
classifier for the prepared four-class branch, evaluated with
specimen-disjoint, class-sensitive metrics. Beyond that, materially
stronger structural or remaining-life claims will ultimately require new
experimental evidence — additional independent campaigns, decoupled
design factors, more structural-labelled specimens, repeated structural
measurements, and observed lifetime data — rather than further model
refinement alone.
